\subsection{Three-generator closure and rank-four trace cubics}
\label{sec:gradient-square-rank-two}

For the closure and realizations in
\cref{prop:cubic-singular-normalization,prop:cubic-gradient-classification},
we determine area minimization, trace models, and the sharp stability
constant. The degree-seven polynomial $G=SF$ complements the common
degree-six potential $F\sqrt{RS}$.

\subsubsection{The tangent obstruction and the existence theorem}

The following closure criterion retains its displayed generality:
for example, the reducible polynomial \(F=xyz\) on \(\R^3\)
satisfies \eqref{eq:cubic-normalized-closure} with \(k=1\), by
direct differentiation. This comparison endpoint is not subject to
the introductory irreducibility convention.

\begin{theorem}[Area minimization from gradient-square closure]
\label{thm:cubic-rank-two-existence}
Let \(F\) be a nonzero homogeneous cubic on \(\R^N\), \(N\geq3\),
and let \(S=|\nabla F|^2\).  Suppose that
\(\R[R,F,S]\) is a Laplacian algebra and that \(F\) has a
nonzero critical point.  Write \(N=6k-3\) as in
\cref{prop:cubic-singular-normalization}.  Then
\begin{equation}\label{eq:cubic-minimality-threshold}
 \partial\llbracket\{F<0\}\rrbracket\text{ is locally mass minimizing}
 \quad\Longleftrightarrow\quad k\geq4.
\end{equation}
Every minimizing member is strictly area minimizing in the sense of
\cref{def:four-variational-notions}, with polynomial potential
\(G=SF\) of degree seven. It is also strictly stable, with the
optimal Hardy constant determined in \cref{thm:sharp}.
\end{theorem}

\begin{proof}
Normalize as in \eqref{eq:cubic-normalized-closure} and suppose
\(k\geq4\). The dictionary in
\cref{prop:cubic-multiplication-dictionary} translates the normalized
gradient identities into the product identities of
\cref{thm:cubic-algebraic-trace-certificate}. Moreover, Bochner's
identity gives
\[
 K=\Delta S=4kR,\qquad
 \tr(L_x^2)=\tfrac12K=2kR\geq8R.
\]
That theorem therefore gives the polynomial subcalibration potential
\(G=SF\), including the required sign and the removability of its
critical set. Its divergence is strictly signed wherever \(F\ne0\).
Since \(G\) is homogeneous and polynomial and both phases are nonempty,
\cref{cor:polynomial-strict-minimality} gives the strict mass gap.

The sharp stability constant is computed in \cref{thm:sharp}.
Harmonicity and \cref{prop:minimizing-irreducibility} also show that
every minimizing member is irreducible.

For the converse, \cref{prop:cubic-gradient-classification} identifies
\(k=1\) with the reducible harmonic cubic \(xyz\), whose phases cannot
minimize perimeter by \cref{prop:minimizing-irreducibility}.
For \(k=2,3\), take a nonzero critical point \(p\).
The Taylor expansion in \cref{prop:cubic-singular-normalization} gives
\[
 r^{-2}F(p+rz)\longrightarrow
 \frac{|p|}{2}\bigl(|u|^2-|v|^2\bigr),
 \qquad z=(u,v,w),\quad u,v\in\R^k.
\]
The transverse quadratic cone is unstable, since its link has
\(|A|^2=2k-2\) and
\[
 \frac{(2k-3)^2}{4}-(2k-2)<0.
\]
By \cref{sec:tangent-instability-transfer}, the Taylor expansion gives
smooth regular graphical convergence. \Cref{lem:tangent-instability-transfer}
therefore makes these cases unstable and excludes area minimization.
\end{proof}

\subsubsection{The rank-four trace realizations}

The trace cubics on $\operatorname{Herm}_4(\mathbb F_d)_0$,
$d=1,2,4$, realize the three exceptional cases in dimensions
$9,15,27$ \cite[Section~6.11.6(iii)]{NadirashviliTkachevVladut14}.
They admit the following uniform calculation.

\begin{corollary}[The rank-four trace cubics]
\label{cor:quaternionic-trace-cubic}
Let \(\mathbb F_d=\R,\mathbb C,\mathbb H\), with real dimension
\(d=1,2,4\), and equip
\[
 V=\operatorname{Herm}_4(\mathbb F_d)_0,\qquad
 \dim_{\R}V=3+6d,
\]
with \(\langle A,B\rangle=\operatorname{Re}\tr(AB)\).
For \(F(A)=\tfrac13\operatorname{Re}\tr(A^3)\),
\(R=\tr(A^2)\), and \(S=|\nabla F|^2\), one has
\[
 \cA_R(F)=\R[R,F,S],\qquad
 \Delta F=0,\quad \Delta S=4(d+1)R,
\]
\[
 \Gamma(F,S)=4RF,\qquad \Gamma(S,S)=4RS+12F^2.
\]
The generators are algebraically independent. The oriented zero cone
is area minimizing if and only if \(d=4\); in that case \(G=SF\)
proves strict area minimality. The common potential \(F\sqrt{RS}\)
also gives a subcalibration in this case.
\end{corollary}

\begin{proof}
Projection onto the trace-free subspace gives
\[
\begin{gathered}
 \nabla F=A^2-\tfrac R4\Id,\qquad
 S=\tr(A^4)-\tfrac{R^2}{4},\\
 (\Hess F)_A(B)=AB+BA-\tfrac12\operatorname{Re}\tr(AB)\Id.
\end{gathered}
\]
Write \(p_j=\sum_i\lambda_i^j\) after diagonalizing \(A\), where
\(p_1=0\). If \(P\) is the orthogonal projection onto the
sum-zero diagonal vectors, the diagonal Hessian is
\(2P\operatorname{diag}(\lambda_i)P\), with trace zero and
squared Hilbert--Schmidt norm \(2R\). Each off-diagonal entry
contributes \(d\) eigenvalues \(\lambda_i+\lambda_j\). Thus
\[
 \Delta F=0,\qquad
 |\Hess F|^2=2R+d\sum_{i<j}(\lambda_i+\lambda_j)^2
             =2(d+1)R,
\]
and \(\Delta S=4(d+1)R\). Also,
\(\nabla S=4A^3-RA-p_3\Id\). Newton's identities give
\[
 p_5=\tfrac56Rp_3,\qquad
 p_6=\tfrac34Rp_4-\tfrac18R^3+\tfrac13p_3^2.
\]
Substitution proves the two gradient identities. Together with
Euler's identities, \cref{prop:MR-generator-test} gives Laplacian
closure. Since \(S=\Gamma(F,F)\) belongs to \(\cA_R(F)\),
this is the smallest algebra containing \(R,F\).
At distinct eigenvalues, the Vandermonde determinant shows that
\(p_2,p_3,p_4\) have independent differentials on \(p_1=0\);
hence \(R,F,S\) are algebraically independent.

The gradient formula and the transverse model in
\cref{prop:cubic-singular-normalization} identify the singular set as
\[
 \Sing\{F=0\}
 =\{A\in V:A^2=\tfrac14\tr(A^2)\Id\}.
\]
Its nonzero matrices have eigenvalues \((\lambda,\lambda,-\lambda,-\lambda)\),
with \(\lambda>0\). In particular,
\(\operatorname{diag}(1,1,-1,-1)\) is a nonzero critical point.
Since \(\Delta S=4kR\) with \(k=d+1\),
\cref{thm:cubic-rank-two-existence} proves all the minimizing and
nonminimizing assertions. For the common potential, put \(D=F^2\).
The displayed identities imply
\[
 \Gamma(F,D)=2SF,\qquad\Gamma(S,D)=8RD,\qquad
 \Gamma(D,D)=4SD.
\]
For \(d=4\), \(N=27\) and \(\tr L_x^2=10R\), so
\cref{cor:cubic-unified-trace} applies.
\end{proof}

\Needspace{12\baselineskip}
\subsubsection{Exact link geometry and the sharp Hardy constant}
On the sphere write \(f=F|_{\Sph^{N-1}}\) and
\(s=S|_{\Sph^{N-1}}\).
Let \(M=Z^*\cap \Sph^{N-1}\). Since \(Z^*\) is conical and radial
vectors are transverse to the sphere, \(M\) is a compact smooth
submanifold of dimension \(4k-4\). The zero link
\(\Sigma=\{F=0\}\cap \Sph^{N-1}\) has regular part of dimension \(6k-5\).
Its transverse dimension along \(M\) is thus \(2k-1\).
If \(\rho\) denotes spherical distance to \(M\), then, near \(M\),
\begin{equation}\label{eq:tube}
 s\asymp\rho^2,\qquad
 \cH^{N-2}\bigl(\Sigma_{\reg}\cap\{\rho<\eps\}\bigr)
 \leq C\eps^{2k-1}.
\end{equation}
The proof of these estimates is given in \cref{sec:cubic-local-geometry}.

\paragraph{Two geometric identities.}
With \(\nu=\nabla F/\sqrt S\), the curvature identity proved in
\cref{prop:harmonic-power-Jacobi} is
\(|A_C|^2=(\Delta S)/(2S)-|\nabla S|^2/(2S^2)\).
Substituting \(\Delta S=4kR\) and \(|\nabla S|^2=4RS\)
on the regular cone gives
\begin{equation}\label{eq:Aexact}
 |A_C|^2=\frac{2(k-1)R}{S},\qquad
 |A_\Sigma|^2=\frac{2(k-1)}s.
\end{equation}

The spherical cross term \(\Gamma_{\Sph}(f,s)=4f(1-3s)\) vanishes on
the zero link. Therefore
\begin{equation}\label{eq:linksgrad}
 |\nabla_\Sigma s|^2=4s-16s^2.
\end{equation}
In particular \(0<s\leq1/4\) on \(\Sigma_{\mathrm{reg}}\), improving
the global bound \(s\leq1/3\) in
\eqref{eq:cubic-gradient-only-S-bound}.

To compute the link Laplacian, differentiate the closure table:
\begin{align*}
 \Gamma(F,\Gamma(F,S))&=4RS+24F^2,\\
 \Hess F(\nabla F,\nabla S)&=\tfrac12|\nabla S|^2=2RS+6F^2.
\end{align*}
Their difference is \(\Hess S(\nabla F,\nabla F)\); hence
\(\Hess S(\nu,\nu)=2R\) on the cone. Since \(S\) has degree four,
its spherical normal Hessian at \(r=1\) equals \(2-4s\).
Subtract this from \(\Delta_{\Sph}s=4k-4(N+2)s\). The link is minimal,
so there is no mean-curvature term, and
\begin{equation}\label{eq:linksLap}
 \Delta_\Sigma s=4k-2-4(N+1)s
 =4k-2-(24k-8)s.
\end{equation}

\paragraph{A positive eigenfunction and its domain.}
Let \(J_\Sigma=\Delta_\Sigma+|A_\Sigma|^2\).
For any real \(p\), equations \eqref{eq:Aexact}--\eqref{eq:linksLap} give
\begin{equation}\label{eq:powerJ}
 J_\Sigma(s^{-p})
 =\left\{\frac{4p^2+(6-4k)p+2k-2}{s}
           +4p(6k-6-4p)\right\}s^{-p}.
\end{equation}
For \(k\geq4\), define
\begin{equation}\label{eq:p}
 \delta_k=\sqrt{4k^2-20k+17}>0,\qquad
 p_k=\frac{2k-3-\delta_k}{4}>0,\qquad
 B_k=4p_k(6k-6-4p_k).
\end{equation}
The coefficient of \(s^{-1}\) in \eqref{eq:powerJ} vanishes, so
\begin{equation}\label{eq:eig}
 \phi=s^{-p_k}>0,\qquad J_\Sigma\phi=B_k\phi.
\end{equation}

\begin{theorem}[Optimal Hardy stability constant]\label{thm:sharp}
Under \eqref{eq:cubic-normalized-closure}, with a nonvertex critical point and \(k\geq4\),
the Friedrichs bottom of \(-J_\Sigma\), initially defined on
\(C_c^\infty(\Sigma_{\mathrm{reg}})\), is \(-B_k\).
The optimal constant in
\[
 Q_C(u)\geq H_k\int_Cr^{-2}u^2,
 \qquad u\in C_c^\infty(C_{\mathrm{reg}}\setminus\{0\}),
\]
is
\begin{equation}\label{eq:sharpH}
 \boxed{H_k=\bigl(k+\sqrt{4k^2-20k+17}\bigr)^2.}
\end{equation}
In particular, for \(k=5\), \(N=27\), \(H_5=42+10\sqrt{17}\).
\end{theorem}

\begin{proof}
For every compactly supported smooth \(v\) on the regular link,
the positive-solution identity applied to \eqref{eq:eig} gives
\begin{equation}\label{eq:linkground}
 \int_\Sigma\bigl(|\nabla v|^2-|A_\Sigma|^2v^2\bigr)
 =\int_\Sigma\phi^2\left|\nabla(v/\phi)\right|^2-B_k\int_\Sigma v^2.
\end{equation}
The bottom is therefore at least \(-B_k\).

By \eqref{eq:tube}, \(\phi\asymp\rho^{-2p_k}\), and \(2k-3-4p_k=\delta_k>0\).
Thus \(\phi\in L^2(\Sigma)\), with finite gradient and potential
energy near \(M\). Choose smooth cutoffs \(\eta_\epsilon\), zero
when \(\rho<\epsilon\), one when \(\rho>2\epsilon\), with
\(|\nabla\eta_\epsilon|\leq C/\epsilon\). The local volume bound gives
\[
 \int_\Sigma\phi^2|\nabla\eta_\epsilon|^2
 =O(\epsilon^{2k-3-4p_k})=O(\epsilon^{\delta_k})\longrightarrow0.
\]
Use \(v=\eta_\epsilon\phi\) in \eqref{eq:linkground}. Its Rayleigh
quotient tends to \(-B_k\), proving the claimed bottom.
The same estimates give convergence to \(\phi\) in the Friedrichs
form norm. Testing \eqref{eq:eig} first against compactly supported
smooth functions and then using form density shows that \(\phi\)
belongs to the Friedrichs operator domain, with eigenvalue \(-B_k\)
for \(-J_\Sigma\).

The cone dimension is \(n=N-1=6k-4\).  By
\eqref{eq:exact-Hardy-gap}, the optimal cone constant is
\[
 \frac{(n-2)^2}{4}+\inf\operatorname{spec}(-J_\Sigma)
 =9(k-1)^2-B_k.
\]
Finally,
\[
 9(k-1)^2-B_k
 =\bigl(3(k-1)-4p_k\bigr)^2=(k+\delta_k)^2,
\]
which proves \eqref{eq:sharpH}.
\end{proof}

\begin{remark}
For the larger root in \eqref{eq:p},
\(2k-3-4p=-\delta_k\), so the cutoff energy does not tend to zero.
The radicand \(\delta_k^2\) equals four times the Hardy stability
constant of the transverse balanced quadratic cone:
\[
 \delta_k^2=4\left[\frac{(2k-3)^2}{4}-(2k-2)\right].
\]
\end{remark}
